\section{System model}
\label{sec:model}

\subsection{Community, and what a boundary is}

A \emph{community} is a region of the stake graph, not an object: nothing in the ledger
names one. There are accounts, underwriters and stakes, and a community is a set of
accounts reachable from a set of underwriters. The boundary between two of them is
\textbf{the absence of a path}.

\textbf{One ledger therefore carries as many communities as it likes.} Two co-operatives
sharing one total order stay separate for exactly as long as nobody trades across --- no
stake edge joins them, so no flow crosses --- and join the moment somebody does, with no
transition and no ceremony. Measured: two co-ops on one chain, each backed by its own
underwriter at $2500$, draw $2500$ apiece and $5000$ together; one settled trade across,
and the first underwriter reaches the second co-op.

The boundary that matters is drawn around the \emph{total order}: one replicated ledger,
one state, one validator set, every quantity in the model computed within it, and
validators that are institutions the ledger's communities themselves operate. Sharing a
ledger is \emph{safer} than not sharing one: a cut is taken over whatever set of accounts
one names, so a coalition spanning two co-operatives on one ledger is measured as one set,
while nothing measures a coalition spanning two deployments --- the scope caveat
\S\ref{sec:security}'s guarantees carry.

\subsection{One ledger}
\label{sec:oneledger}

This design contemplates one ledger. A community does not get its own: one order carries
as many regions as its members form, so the case that would call for a second --- two
communities that want to trade --- is the case a shared order already answers.

The rule that holds it is \textbf{nothing outside a ledger's own total order may authorise
a transition inside it}, and every discharge in the alphabet already obeys it, being
authorised by the party who loses if it is wrong: settlement and cure by the creditor's
signature, an arbitral award by a panel both parties named, capped and time-boxed before
the obligation existed.

What the rule refuses is an order honouring proofs drawn from another's state root. An
inclusion proof certifies that foreign validators committed a leaf, and for that to be
worth more than a signature the honouring ledger must check it against a root it
\emph{knows} to be theirs, which means tracking a foreign validator set through its
changes; anything weaker is a signature by a party who could have signed an ordinary
settlement instead. So it is a light client or it is nothing. Measured against a community
of six underwriters at $2500$: honouring a foreign declared supply yields $6\times10^{9}$
where a shared order yields exactly zero, honouring a proof as discharge extracts the
community's entire ceiling on \emph{no} local signatures, and a forged discharge recycles
to $1000\times$ the attacker's real backing.

Two findings survive the refusal. A genesis file is as free as a key: the design rests on
identities being free while stakes are not, so a cut drawn over one order's accounts never
measures a coalition that stood up a second. And a forged discharge is cheaper than a
default and repeatable, where a defaulter spends their standing once.

A separate deployment is therefore a separate system, with no credit path to it and no
protocol for one. A member who wants to participate there trades there --- the first trade
seats an account and grants nothing --- and earns standing there as everybody does.

What bounds one ledger bounds the design, stated as a bound rather than escaped by
splitting. \S\ref{sec:implementation} measures the binding cost as the per-block re-audit
rather than the per-acceptance query: $115$\,ms on the epoch-advance block over
$20{,}000$ accounts with $200$ underwriters --- the most expensive block there is ---
against $19$\,ms for one acceptance. No term grows with the underwriter count, because an
obligation's stored hold is a witness the audit verifies rather than a cut it recomputes;
what remains grows with the graph. Carrying substantially more than that on one order is
open, and splitting does not answer it: severing a stake graph destroys every edge
crossing the cut, and two orders becoming one is a founding rather than a migration, so
nothing rejoins them.

\subsection{Time}

Time advances in \emph{epochs} of $\kConstEpochSecs$ seconds. An epoch number is
absolute --- the timestamp divided by the constant, identical on every ledger --- and a
block carries a timestamp from which the ledger derives its target epoch and closes
epochs one at a time. What runs \emph{per epoch closed} is what a missed epoch would
have to be replayed for: stake decay, the free-allowance and saturation counters, and
replay-cache pruning.

What runs \emph{once} at the end of the advance is the sweep: the permissionless cranks
that default a matured obligation and forfeit the bonds of a member who has exhausted the
write gate for $\kConstBondForfeitEpochs$ consecutive epochs, followed by bond release.
Both are comparisons against the epoch the block ends at, so one pass sees exactly what a
pass per epoch would have seen, and a fresh chain climbing thousands of epochs toward
wall-clock in one block does not walk the contract book thousands of times. The order of
the two is load-bearing: releasing bonds before the forfeiture crank looked would hand
back the encumbrance the sanction exists to take. Running the default crank here, instead
of leaving it to whoever cares, is what makes \S\ref{sec:recourse}'s substitution
automatic: no obligation is ever due and uncranked, so a creditor need not notice their
own loss to be made whole.

Epoch boundaries are the only scheduled events: no continuous process, no timer and no
off-chain job. Every state change is a transition in the total order or a deterministic
consequence of advancing the epoch counter.

\subsection{State}

The replicated state is, with every amount in it an integer number of minor units
(\S\ref{sec:implementation}):

\begin{itemize}[leftmargin=1.4em]
  \item \textbf{Accounts.} A set of keys, optional guardian configuration for key
        recovery, cached outstanding debt, how much of it is currently in default, the
        two support-cascade declarations of \S\ref{sec:medium}, operation-bond
        bookkeeping, and two advisory debt-velocity counters that only a client reads.
        There is no score, no tier, no attestation and no sponsor.

        An account also carries a \emph{status}, which is not an admission ladder and
        along which nothing is promoted: \textbf{active} is the resting state of every
        account that exists, and the other two are a governance sanction
        (\textbf{suspended}, \S\ref{sec:governance}) and a voluntary departure
        (\textbf{exited}, refused while the member owes anything, is in default, holds an
        unreleased bond, still has a supply declared, or is the last validator). Neither
        is consulted by Definition~\ref{def:capacity}, which is what \S\ref{sec:standing}
        means by an account having no status.
  \item \textbf{The stake graph.} $\stake(c,d)$ in minor units, \emph{directed}, keyed by
        the ordered pair creditor-to-debtor (Definition~\ref{def:stake}); held
        symmetrically it would say that extending credit raises the lender's own limit,
        which the model does not mean.
  \item \textbf{Reservations.} $\res(c,d)$, keyed identically, and the committed total per
        underwriter, together recording the flow held against outstanding insured credit
        (Definition~\ref{def:reservation}). Both halves are state because both are charged
        (Remark~\ref{rem:charge-supply}).
  \item \textbf{Underwriters.} Members who have declared a supply
        (Definition~\ref{def:underwriter}), with the amount each declared, every unit of
        which arrived through a ceremony rather than against standing the community
        itself conferred (\S\ref{sec:governance}). Genesis seeds the first ones;
        thereafter any member may join or leave the set.
  \item \textbf{Contracts.} Debtor, creditor, outstanding and original amount, maturity,
        the epoch of acceptance, status, whether the obligation is insured, \emph{exactly
        the flow it holds} --- which arcs and how much on each
        (Remark~\ref{rem:charge-supply}) --- and consented arbitration terms if any.
  \item \textbf{Governance.} Proposals, governed constants, validator powers.
  \item \textbf{Chain identity.} Chain id, root salt, replay cache.
\end{itemize}

There is no reputation structure --- no per-relationship evidence, no decayed counter, no
reliability estimate, no volume envelope: the stake graph answers directly the question
all of those approximated.

\subsection{Observables}

A member's client can compute, from public state, its own capacity and who supplies it,
its outstanding obligations and their maturities, whether a counterparty's proposed
obligation would be insured, and community-level totals. Precise amounts may be sealed by
charter policy to the parties and validators, with others seeing bucketed magnitudes
(\S\ref{sec:implementation}). Capacity is a pure function of committed state, so any
observer holding the ledger computes the value any validator does: checkable, not
announced.

\subsection{Notation}

$\uw$ is the set of underwriters and $\supply(u)$ the supply $u$ declared.
$\stake(c,d)$ is the directed stake of creditor $c$ on debtor $d$ and $\res(c,d)$ the
reservation against it, with $\free(c,d)=\stake(c,d)-\res(c,d)$. $\Capa(T)$ is the
capacity of an account set $T$, and $\Capa(x)$ abbreviates $\Capa(\{x\})$. $\partial S$
is the set of stakes with exactly one endpoint in $S$. $\vbase$ is the reference
denomination, genesis $\kConstBaseCapacity$.
